%& -shell-escape
% !TeX root = thesis-abstract.tex
% !TeX encoding = UTF-8
% !TeX program = XeLaTeX

\documentclass{.local/lib/classes/ndvr-super-article-standalone}

\begin{document}

    Работать с~физическим представлением видео неудобно.
    Для практической реализации поиска
    используются \textit{сигнатуры} видео. 
    Сигнатура видео — компактное представление объекта,
    обеспечивающее его идентификацию.

    \begin{minipage}{\textwidth}
        Поиск по~видео делится на~три основных этапа:
        \begin{itemize}
            \item создание сигнатур;
            \item управление сигнатурами (создание индекса);
            \item поиск сигнатур. 
        \end{itemize}
    \end{minipage}

    В~некоторых приложениях
    существуют дополнительные этапы анализа, 
    необходимые для~более
    точного определения положения дубликатов.
    На рисунке \ref{fig:video-search-dfd} приведена диаграмма
    потоков данных для поиска по видео.

    % \begin{figure}[H]
    %     \includegraphics[width=\textwidth]
    %         {.local/vec/out/pdf/video-model-dfd-state-x.pdf}
    %     \caption{DFD-диаграмма поиска по видео 
    %         в нотации Йордона — Де Марко.}
    %     \label{fig:video-search-dfd-x}
    % \end{figure}

    Чаще всего методы анализа видео различаются
    способом формирования сигнатур 
    из низкоуровневых характеристик\footnote{
        Liu, J., Huang, Z., Cai, H., Shen, H. T., Ngo, C. W., a
        nd Wang, W. 2013. Near-duplicate video retrieval:
        Current research and future trends. 
        ACM Comput. Surv. 45, 4, Article 44 (August 2013), 
        23 pages.
        DOI: http://dx.doi.org/10.1145/2501654.2501658
    }.
    Алгоритмы выделения характеристик кадров получены
    в~результате исследований 
    в~области анализа изображений.

    В работе рассмотрены четыре основных типа сигнатур: 
    \begin{itemize}%[nosep]
        \item {
            локальные сигнатуры уровня кадра —
            выражают локальную информацию, 
            содержащуюся в~кадре;
        }
        \item {
            глобальные сигнатуры уровня кадра — 
            выражают одной сигнатурой целый кадр;
        }
        \item {
            глобальные сигнатуры уровня видео — 
            выражает одной сигнатурой все видео целиком;
        }
        \item {
            пространственно-временные~сигнатуры — 
            выражают видео целиком, 
            используя информацию 
            об~изменении кадров во~времени.
        }
    \end{itemize}


    Многие методы используют только пространственные 
    характеристики видео\footnote{
        W. Hu, N. Xie, L. Li, X. Zeng and S. Maybank, 
        A Survey on Visual Content-Based Video Indexing and Retrieval, 
        in IEEE Transactions on Systems, Man, and Cybernetics, 
        Part C (Applications and Reviews), 
        vol. 41, no. 6, pp. 797-819, Nov. 2011.
        doi: 10.1109/TSMCC.2011.2109710
    }.
    Изменение кадров во времени никак не учитывается.
    При этом, при построении сигнатуры в неё закладывают 
    второстепенные характеристики кадров, 
    которые мешают поиску дубликатов.
    Именно поэтому, пространственно-временные~сигнатуры 
    представляют особый интерес.

    Видео — это не просто последовательность кадров, 
    а последовательность фактов или событий,
    развивающихся во времени.
    Темпоральная модель видео представлена 
    на рисунке \ref{fig:temporal-uml}.

    Работа с событиями в видео подразумевает разбиение
    видео на отрезки и~поиск событий внутри этих отрезков.
    Однако сам метод разбиения в~значительной степени влияет 
    на~формирование итоговой сигнатуры.

   \begin{figure}[p]
        \includegraphics[width=\textwidth]
            {.local/vec/out/pdf/video-model-uml-temporal.pdf}
        \caption{
            Темпоральная модель видео,\\
            выраженная диаграммой классов UML 2.0%
        }
        \label{fig:temporal-uml}
    \end{figure}

\end{document}
